% The seven edges forced by the textbook's own examples (a subgraph of its Fig. 1.2.5).
\setlength{\unitlength}{1pt}
\begin{picture}(135,90)(-25,-10)
  \put(0,35){\circle*{6}} \put(40,70){\circle*{6}} \put(40,0){\circle*{6}}
  \put(80,70){\circle*{6}} \put(80,0){\circle*{6}}
  \put(0,35){\line(4,3){40}}   % v1 v2
  \put(0,35){\line(4,-3){40}}  % v1 v3
  \put(40,70){\line(0,-1){70}} % v2 v3
  \put(40,70){\line(1,0){40}}  % v2 v4
  \put(40,0){\line(2,3){40}}   % v3 v4
  \put(40,0){\line(1,0){40}}   % v3 v5
  \put(80,70){\line(0,-1){70}} % v4 v5
  \put(-20,32){$v_1$} \put(33,76){$v_2$} \put(33,-12){$v_3$}
  \put(84,70){$v_4$} \put(84,-3){$v_5$}
\end{picture}
